% Curs 12 - Pagina 8 (3-I-1995).
% Fragment pentru % Curs 12 - Pagina 8 (3-I-1995).
% Fragment pentru % Curs 12 - Pagina 8 (3-I-1995).
% Fragment pentru % Curs 12 - Pagina 8 (3-I-1995).
% Fragment pentru \input{courses/pages/curs12-pagina8-ro}.
% Fara preambul sau salt de pagina fortat.
% NOTE EDITORIALE:
% - Pagina continua constructia grupului liber FG[A] de la pagina 7.
% - In enuntul proprietatii universale, tinta aplicatiei zeta seamana
%   cu G^* in manuscris, desi grupul din diagrama este G. Se pastreaza
%   citirea aparenta G^*, fara corectie tacita.
% - Textul si diagrama a doua sunt foarte comprimate. Zeta_1 este
%   prelungirea functiei zeta la alfabetul A reunit cu A'.
% - Pe linia f o zeta_{A unite A'} = zeta_1, litera poate semana cu
%   f sau o alta litera; lectura f este determinata de continuarea
%   demonstratiei si ramane provizorie.
% - Manuscrisul noteaza (xx',1) in R si verifica f(xx')=f(1).
% - Ultimele doua randuri pun problema cuvintelor sub forma unei
%   intrebari despre un algoritm finit; raspunsul foarte palid se
%   citeste probabil "nu totdeauna".

\subsection*{Curs 12 - Pagina 8}

\noindent\fbox{T}
\begin{enumerate}
\item[(1)] $\mathrm{FG}[A]$ este grup.
\item[(2)] Aplica\c{t}ia $\zeta_A:A\longrightarrow\mathrm{FG}[A]$
este injectiv\u{a}.
\end{enumerate}

\medskip
\noindent $(\forall)\,(G,\cdot,e)$ grup,
$(\forall)\,\zeta:A\longrightarrow G^*$,
\[
\exists!\,g:\mathrm{FG}[A]\longrightarrow G
\quad\text{a.\^{i}.}\quad g\circ\zeta_A=\zeta.
\]

\[
\begin{tikzcd}[column sep=large, row sep=large]
A \arrow[r,"\zeta_A"] \arrow[dr,"\zeta"']
  & \mathrm{FG}[A] \arrow[d,dashed,"g"]\\
& G
\end{tikzcd}
\]

\medskip
\noindent $A\subseteq A\cup A'$,
\[
A\cup A'\xrightarrow{\ \zeta_{A\cup A'}\ }(A\cup A')^*
\xrightarrow{\ \rho\ }
\frac{(A\cup A')^*}{R^*}
=\langle\overline{A}\cup\overline{A'}\rangle.
\]

\[
\begin{tikzcd}[column sep=large, row sep=large]
A\cup A' \arrow[r,"\zeta_{A\cup A'}"] \arrow[dr,"\zeta_1"']
  & (A\cup A')^* \arrow[r,"\rho"] \arrow[d,dashed,"f"]
  & \mathrm{FG}[A] \arrow[dl,dashed,"g"]\\
& G &
\end{tikzcd}
\]

\noindent
\[
\zeta_1(x)=\zeta(x)\quad(\forall)\,x\in A,
\qquad
\zeta_1(x')=\zeta(x)^{-1}\quad(\forall)\,x'\in A'.
\]

\noindent $f\circ\zeta_{A\cup A'}=\zeta_1$,
\qquad $R^*\subseteq R_f$.

\medskip
\noindent $(xx',1)\in R\subseteq R_f$.
\[
f(xx')=f(x)f(x')=e=f(1),
\]
\noindent unde $f(x)=\zeta_1(x)=\zeta(x)$ \c{s}i
$f(x')=\zeta_1(x')=\zeta(x)^{-1}$.

\bigskip
\noindent\underline{\textit{``Word problem''} (problema cuvintelor)}
\[
R\subseteq A^*\times A^*,\qquad
R^*\subseteq A^*\times A^*,\qquad
u,v\in A^*,\quad u\,R^*\,v.
\]

\noindent $\exists$ algoritm a.\^{i}.\ ob\c{t}inem $v$ din $u$
\^{i}n num\u{a}r finit de pa\c{s}i?

\medskip
\noindent\textit{Nu totdeauna.}
.
% Fara preambul sau salt de pagina fortat.
% NOTE EDITORIALE:
% - Pagina continua constructia grupului liber FG[A] de la pagina 7.
% - In enuntul proprietatii universale, tinta aplicatiei zeta seamana
%   cu G^* in manuscris, desi grupul din diagrama este G. Se pastreaza
%   citirea aparenta G^*, fara corectie tacita.
% - Textul si diagrama a doua sunt foarte comprimate. Zeta_1 este
%   prelungirea functiei zeta la alfabetul A reunit cu A'.
% - Pe linia f o zeta_{A unite A'} = zeta_1, litera poate semana cu
%   f sau o alta litera; lectura f este determinata de continuarea
%   demonstratiei si ramane provizorie.
% - Manuscrisul noteaza (xx',1) in R si verifica f(xx')=f(1).
% - Ultimele doua randuri pun problema cuvintelor sub forma unei
%   intrebari despre un algoritm finit; raspunsul foarte palid se
%   citeste probabil "nu totdeauna".

\subsection*{Curs 12 - Pagina 8}

\noindent\fbox{T}
\begin{enumerate}
\item[(1)] $\mathrm{FG}[A]$ este grup.
\item[(2)] Aplica\c{t}ia $\zeta_A:A\longrightarrow\mathrm{FG}[A]$
este injectiv\u{a}.
\end{enumerate}

\medskip
\noindent $(\forall)\,(G,\cdot,e)$ grup,
$(\forall)\,\zeta:A\longrightarrow G^*$,
\[
\exists!\,g:\mathrm{FG}[A]\longrightarrow G
\quad\text{a.\^{i}.}\quad g\circ\zeta_A=\zeta.
\]

\[
\begin{tikzcd}[column sep=large, row sep=large]
A \arrow[r,"\zeta_A"] \arrow[dr,"\zeta"']
  & \mathrm{FG}[A] \arrow[d,dashed,"g"]\\
& G
\end{tikzcd}
\]

\medskip
\noindent $A\subseteq A\cup A'$,
\[
A\cup A'\xrightarrow{\ \zeta_{A\cup A'}\ }(A\cup A')^*
\xrightarrow{\ \rho\ }
\frac{(A\cup A')^*}{R^*}
=\langle\overline{A}\cup\overline{A'}\rangle.
\]

\[
\begin{tikzcd}[column sep=large, row sep=large]
A\cup A' \arrow[r,"\zeta_{A\cup A'}"] \arrow[dr,"\zeta_1"']
  & (A\cup A')^* \arrow[r,"\rho"] \arrow[d,dashed,"f"]
  & \mathrm{FG}[A] \arrow[dl,dashed,"g"]\\
& G &
\end{tikzcd}
\]

\noindent
\[
\zeta_1(x)=\zeta(x)\quad(\forall)\,x\in A,
\qquad
\zeta_1(x')=\zeta(x)^{-1}\quad(\forall)\,x'\in A'.
\]

\noindent $f\circ\zeta_{A\cup A'}=\zeta_1$,
\qquad $R^*\subseteq R_f$.

\medskip
\noindent $(xx',1)\in R\subseteq R_f$.
\[
f(xx')=f(x)f(x')=e=f(1),
\]
\noindent unde $f(x)=\zeta_1(x)=\zeta(x)$ \c{s}i
$f(x')=\zeta_1(x')=\zeta(x)^{-1}$.

\bigskip
\noindent\underline{\textit{``Word problem''} (problema cuvintelor)}
\[
R\subseteq A^*\times A^*,\qquad
R^*\subseteq A^*\times A^*,\qquad
u,v\in A^*,\quad u\,R^*\,v.
\]

\noindent $\exists$ algoritm a.\^{i}.\ ob\c{t}inem $v$ din $u$
\^{i}n num\u{a}r finit de pa\c{s}i?

\medskip
\noindent\textit{Nu totdeauna.}
.
% Fara preambul sau salt de pagina fortat.
% NOTE EDITORIALE:
% - Pagina continua constructia grupului liber FG[A] de la pagina 7.
% - In enuntul proprietatii universale, tinta aplicatiei zeta seamana
%   cu G^* in manuscris, desi grupul din diagrama este G. Se pastreaza
%   citirea aparenta G^*, fara corectie tacita.
% - Textul si diagrama a doua sunt foarte comprimate. Zeta_1 este
%   prelungirea functiei zeta la alfabetul A reunit cu A'.
% - Pe linia f o zeta_{A unite A'} = zeta_1, litera poate semana cu
%   f sau o alta litera; lectura f este determinata de continuarea
%   demonstratiei si ramane provizorie.
% - Manuscrisul noteaza (xx',1) in R si verifica f(xx')=f(1).
% - Ultimele doua randuri pun problema cuvintelor sub forma unei
%   intrebari despre un algoritm finit; raspunsul foarte palid se
%   citeste probabil "nu totdeauna".

\subsection*{Curs 12 - Pagina 8}

\noindent\fbox{T}
\begin{enumerate}
\item[(1)] $\mathrm{FG}[A]$ este grup.
\item[(2)] Aplica\c{t}ia $\zeta_A:A\longrightarrow\mathrm{FG}[A]$
este injectiv\u{a}.
\end{enumerate}

\medskip
\noindent $(\forall)\,(G,\cdot,e)$ grup,
$(\forall)\,\zeta:A\longrightarrow G^*$,
\[
\exists!\,g:\mathrm{FG}[A]\longrightarrow G
\quad\text{a.\^{i}.}\quad g\circ\zeta_A=\zeta.
\]

\[
\begin{tikzcd}[column sep=large, row sep=large]
A \arrow[r,"\zeta_A"] \arrow[dr,"\zeta"']
  & \mathrm{FG}[A] \arrow[d,dashed,"g"]\\
& G
\end{tikzcd}
\]

\medskip
\noindent $A\subseteq A\cup A'$,
\[
A\cup A'\xrightarrow{\ \zeta_{A\cup A'}\ }(A\cup A')^*
\xrightarrow{\ \rho\ }
\frac{(A\cup A')^*}{R^*}
=\langle\overline{A}\cup\overline{A'}\rangle.
\]

\[
\begin{tikzcd}[column sep=large, row sep=large]
A\cup A' \arrow[r,"\zeta_{A\cup A'}"] \arrow[dr,"\zeta_1"']
  & (A\cup A')^* \arrow[r,"\rho"] \arrow[d,dashed,"f"]
  & \mathrm{FG}[A] \arrow[dl,dashed,"g"]\\
& G &
\end{tikzcd}
\]

\noindent
\[
\zeta_1(x)=\zeta(x)\quad(\forall)\,x\in A,
\qquad
\zeta_1(x')=\zeta(x)^{-1}\quad(\forall)\,x'\in A'.
\]

\noindent $f\circ\zeta_{A\cup A'}=\zeta_1$,
\qquad $R^*\subseteq R_f$.

\medskip
\noindent $(xx',1)\in R\subseteq R_f$.
\[
f(xx')=f(x)f(x')=e=f(1),
\]
\noindent unde $f(x)=\zeta_1(x)=\zeta(x)$ \c{s}i
$f(x')=\zeta_1(x')=\zeta(x)^{-1}$.

\bigskip
\noindent\underline{\textit{``Word problem''} (problema cuvintelor)}
\[
R\subseteq A^*\times A^*,\qquad
R^*\subseteq A^*\times A^*,\qquad
u,v\in A^*,\quad u\,R^*\,v.
\]

\noindent $\exists$ algoritm a.\^{i}.\ ob\c{t}inem $v$ din $u$
\^{i}n num\u{a}r finit de pa\c{s}i?

\medskip
\noindent\textit{Nu totdeauna.}
.
% Fara preambul sau salt de pagina fortat.
% NOTE EDITORIALE:
% - Pagina continua constructia grupului liber FG[A] de la pagina 7.
% - In enuntul proprietatii universale, tinta aplicatiei zeta seamana
%   cu G^* in manuscris, desi grupul din diagrama este G. Se pastreaza
%   citirea aparenta G^*, fara corectie tacita.
% - Textul si diagrama a doua sunt foarte comprimate. Zeta_1 este
%   prelungirea functiei zeta la alfabetul A reunit cu A'.
% - Pe linia f o zeta_{A unite A'} = zeta_1, litera poate semana cu
%   f sau o alta litera; lectura f este determinata de continuarea
%   demonstratiei si ramane provizorie.
% - Manuscrisul noteaza (xx',1) in R si verifica f(xx')=f(1).
% - Ultimele doua randuri pun problema cuvintelor sub forma unei
%   intrebari despre un algoritm finit; raspunsul foarte palid se
%   citeste probabil "nu totdeauna".

\subsection*{Curs 12 - Pagina 8}

\noindent\fbox{T}
\begin{enumerate}
\item[(1)] $\mathrm{FG}[A]$ este grup.
\item[(2)] Aplica\c{t}ia $\zeta_A:A\longrightarrow\mathrm{FG}[A]$
este injectiv\u{a}.
\end{enumerate}

\medskip
\noindent $(\forall)\,(G,\cdot,e)$ grup,
$(\forall)\,\zeta:A\longrightarrow G^*$,
\[
\exists!\,g:\mathrm{FG}[A]\longrightarrow G
\quad\text{a.\^{i}.}\quad g\circ\zeta_A=\zeta.
\]

\[
\begin{tikzcd}[column sep=large, row sep=large]
A \arrow[r,"\zeta_A"] \arrow[dr,"\zeta"']
  & \mathrm{FG}[A] \arrow[d,dashed,"g"]\\
& G
\end{tikzcd}
\]

\medskip
\noindent $A\subseteq A\cup A'$,
\[
A\cup A'\xrightarrow{\ \zeta_{A\cup A'}\ }(A\cup A')^*
\xrightarrow{\ \rho\ }
\frac{(A\cup A')^*}{R^*}
=\langle\overline{A}\cup\overline{A'}\rangle.
\]

\[
\begin{tikzcd}[column sep=large, row sep=large]
A\cup A' \arrow[r,"\zeta_{A\cup A'}"] \arrow[dr,"\zeta_1"']
  & (A\cup A')^* \arrow[r,"\rho"] \arrow[d,dashed,"f"]
  & \mathrm{FG}[A] \arrow[dl,dashed,"g"]\\
& G &
\end{tikzcd}
\]

\noindent
\[
\zeta_1(x)=\zeta(x)\quad(\forall)\,x\in A,
\qquad
\zeta_1(x')=\zeta(x)^{-1}\quad(\forall)\,x'\in A'.
\]

\noindent $f\circ\zeta_{A\cup A'}=\zeta_1$,
\qquad $R^*\subseteq R_f$.

\medskip
\noindent $(xx',1)\in R\subseteq R_f$.
\[
f(xx')=f(x)f(x')=e=f(1),
\]
\noindent unde $f(x)=\zeta_1(x)=\zeta(x)$ \c{s}i
$f(x')=\zeta_1(x')=\zeta(x)^{-1}$.

\bigskip
\noindent\underline{\textit{``Word problem''} (problema cuvintelor)}
\[
R\subseteq A^*\times A^*,\qquad
R^*\subseteq A^*\times A^*,\qquad
u,v\in A^*,\quad u\,R^*\,v.
\]

\noindent $\exists$ algoritm a.\^{i}.\ ob\c{t}inem $v$ din $u$
\^{i}n num\u{a}r finit de pa\c{s}i?

\medskip
\noindent\textit{Nu totdeauna.}
